\documentclass[10pt]{article}
\usepackage[utf8]{inputenc}
\usepackage[T1]{fontenc}
\usepackage[polish]{babel}
\usepackage{fullpage}
\usepackage{amsmath}
\usepackage{amssymb}
\usepackage{mathtools}
\usepackage{setspace}
\usepackage{booktabs}
\usepackage{array}
\usepackage{hyperref}
\doublespacing

% Pola do uzupełnienia
\newcommand{\uzupelnij}[1]{\textbf{[#1]}}
\newcommand{\autorA}{Jakub Rudnik}
\newcommand{\autorB}{Marcel Kowalczyk}
\newcommand{\repo}{\url{https://github.com/Instytut-Softu-Niepowaznego/Conformal-Risk-Control}}

\begin{document}

\title{Kontrola ryzyka pominięcia etykiet za pomocą Conformal Risk Control\\[0.5em]
\large Deklaracja projektu --- Rozpoznawanie wzorców, studia II stopnia}
\author{\autorA \and \autorB}
\date{\today}

\maketitle

\begin{abstract}

Projekt dotyczy kontroli niepewności i ryzyka predykcji w klasyfikacji wieloetykietowej.
Odtworzono procedurę Conformal Risk Control (CRC) \cite{CRC} w uproszczonym, małym
problemie syntetycznym i zbadano, jak wielkość zbioru kalibracyjnego wpływa na kontrolę
ryzyka pominięcia prawdziwych etykiet oraz na rozmiar zwracanych zbiorów predykcji.
CRC porównano ze stałym progiem prawdopodobieństwa 0{,}5.

\end{abstract}

\pagebreak

%--TEMAT--%
\section{Temat projektu}

\textbf{Temat nr 17:} Kontrola ryzyka pominięcia etykiet za pomocą Conformal Risk Control.

\noindent\textbf{Obszar:} niepewność i kontrola ryzyka predykcji.

\noindent\textbf{Publikacja bazowa:} A.~Angelopoulos, S.~Bates, A.~Fisch, L.~Lei,
T.~Schuster, \emph{Conformal Risk Control}, ICLR 2024 \cite{CRC}.

\noindent\textbf{Pytanie badawcze:} Jak wielkość zbioru kalibracyjnego wpływa na kontrolę
ryzyka pominięcia etykiet i rozmiar zwracanych zbiorów predykcji?

%--OPIS ZADANIA--%
\section{Opis zadania}

\subsection{Podstawy teoretyczne}

Rozważono predykcję wieloetykietową z $K$ etykietami: dla obiektu $X$ zbiór prawdziwych
etykiet $Y \subseteq \{1,\dots,K\}$ jest niepusty, a model zwraca zbiór predykcji
$C_\lambda(X) = \{k : \hat p_k(X) \ge 1-\lambda\}$, gdzie $\hat p_k$ to oszacowane
prawdopodobieństwo etykiety $k$. Stratą jest odsetek pominiętych prawdziwych etykiet:
\[
  \ell\bigl(C_\lambda(X), Y\bigr) = 1 - \frac{|Y \cap C_\lambda(X)|}{|Y|} \in [0, B], \qquad B = 1 .
\]
Strata jest ograniczona i monotoniczna: przy rosnącym $\lambda$ zbiór $C_\lambda$ rośnie,
a strata nie rośnie. Przykładowo dla $Y=\{A,B\}$ predykcja $\{A\}$ daje stratę $1/2$,
a predykcja $\{A,B,C\}$ stratę $0$ --- kosztem dodatkowej, fałszywej etykiety.

Na podstawie $n$ obiektów kalibracyjnych, wymiennych z obiektem testowym, CRC wybiera
\[
  \hat\lambda = \inf\Bigl\{\lambda : \tfrac{n}{n+1}\,\hat R_n(\lambda) + \tfrac{B}{n+1} \le \alpha\Bigr\},
  \qquad \hat R_n(\lambda) = \frac{1}{n}\sum_{i=1}^{n} \ell\bigl(C_\lambda(X_i), Y_i\bigr),
\]
co gwarantuje $\mathbb{E}\bigl[\ell(C_{\hat\lambda}(X_{n+1}), Y_{n+1})\bigr] \le \alpha$.
Gwarancja dotyczy wartości oczekiwanej ryzyka --- nie poprawności każdej pojedynczej
predykcji ani warunkowej kontroli w podgrupach. Sam próg surowego średniego błędu
(bez poprawki $n/(n+1)$ i składnika $B/(n+1)$) nie jest pełną procedurą CRC.

\subsection{Dane i zakres odtworzenia}

\begin{itemize}
  \item Syntetyczny problem z 5 etykietami i 20 cechami, co najmniej jedna prawdziwa
        etykieta na obiekt; kontrolowany generator zapewnia wymienność danych
        kalibracyjnych i testowych.
  \item Rozłączne części: treningowa, kalibracyjna i testowa (kilka tysięcy obiektów
        na zbiór).
  \item Klasyfikator: OneVsRest z regresją logistyczną, trenowany na części niezależnej
        od kalibracji.
  \item Adaptacja procedury CRC do małych danych; docelowe ryzyko $\alpha = 0{,}1$;
        punkt odniesienia: stały próg prawdopodobieństwa $0{,}5$.
  \item Raportowane metryki: średnia strata oraz średnia liczba etykiet w zbiorze
        predykcji.
\end{itemize}

\subsection{Eksperyment własny}

Porównano kalibrację na $n \in \{50, 200, 500\}$ obiektach. Dla każdego rozmiaru
wykonano 30 losowań podziału kalibracja/test i zaraportowano ryzyko, rozmiar zbiorów
predykcji oraz ich zmienność. Parametrów nie dobierano na podstawie zbioru testowego.

\subsection{Metodologia i kontrola poprawności}

\begin{itemize}
  \item Przed końcowym badaniem utrwalono podziały danych, parametry, metryki i ziarna
        losowań.
  \item Zachowano wyniki wszystkich zaplanowanych uruchomień, błędy oraz powody
        ewentualnych brakujących wyników.
  \item Zinterpretowano kontrolę oczekiwanego ryzyka, a nie wymaganie, by każdy skończony
        test dał wynik poniżej $\alpha$; omówiono koszt zwiększania zbioru predykcji.
  \item Udokumentowano zmiany względem publikacji.
\end{itemize}

%--REPOZYTORIUM--%
\section{Repozytorium}

Kod, konfiguracje, surowe wyniki oraz instrukcja uruchomienia:
\begin{center}
\repo
\end{center}

%--HARMONOGRAM--%
\section{Harmonogram i punkty kontrolne}

\begin{center}
\renewcommand{\arraystretch}{1.2}
\begin{tabular}{>{\raggedright}p{0.17\textwidth} >{\raggedright}p{0.53\textwidth} p{0.18\textwidth}}
\toprule
\textbf{Etap} & \textbf{Zakres i materiały} & \textbf{Termin} \\
\midrule
Etap 1: Prezentacja tematu &
Strata, monotoniczność, wymienność; przykład etykiet i poprawki kalibracyjnej;
ustalenie generatora, $\alpha$ i rozłącznych części danych.
Slajdy + protokół ($\sim$1 strona). &
\uzupelnij{data} \\
\midrule
Etap 2: Wstępne wyniki &
Działający klasyfikator, próg stały i pierwsza kalibracja CRC; powiązanie wzoru
poprawki z kodem. Slajdy, uruchamialny kod, pierwsze wyniki, notatka z postępów
(1--2 strony). &
\uzupelnij{data} \\
\midrule
Etap 3: Prezentacja i raport końcowy &
30 losowań dla każdego rozmiaru kalibracji; analiza ryzyka, wielkości predykcji
i założeń gwarancji. Raport 6--8 stron, slajdy, kod, instrukcja uruchomienia,
wersje zależności, surowe wyniki. &
\uzupelnij{data} \\
\bottomrule
\end{tabular}
\end{center}

%--BIBLIOGRAPHY--%
\begin{thebibliography}{99}

\bibitem{CRC} A.~N.~Angelopoulos, S.~Bates, A.~Fisch, L.~Lei, T.~Schuster,
``Conformal Risk Control'', International Conference on Learning Representations
(ICLR), 2024. \\
Publikacja bazowa projektu; wprowadza procedurę CRC kontrolującą oczekiwaną wartość
ograniczonej, monotonicznej straty.

\end{thebibliography}

\end{document}
